%% =============================================================================
%% variaveis.tex — dados da capa, folha de rosto, resumo e abstract
%% Texto em vermelho (\textcolor{red}{...}) = pendente: trocar antes da entrega.
%% =============================================================================

%% Fonte do texto: \sfdefault = Helvetica (equivalente à Arial); \rmdefault = Times
\fonteTexto{\sfdefault}

%% Tipo e nível do trabalho: a classe usa estes valores para escolher a formatação de
%% graduação (reconhece "Tese", "Dissertação", "Trabalho de Conclusão de Curso de Graduação"
%% e "Estágio Supervisionado"). O tipo que aparece no documento é o \documenttype abaixo.
\TipoDeDocumento{Trabalho de Conclusão de Curso de Graduação}
\NivelDeFormacao{Bacharelado}

%% Curso
\program{\textcolor{red}{Bacharelado em Ciência da Computação}}
\programen{Undergradute Program in Computer Science}
\degree{Bacharel}
\degreearea{}

%% Disciplina e natureza do trabalho (folha de rosto)
\goal{Relatório técnico apresentado}
\course{Sistemas Distribuídos}

\makeatletter
\goalfulltext{%
  \begin{Spacing}{0.8}
    \@goal{} como requisito parcial de avaliação da disciplina de
    \@course{} do \@program{} da \@university{}.
  \end{Spacing}
}
\makeatother

%% Autores (um \addauthor por integrante)
\addauthor{\textcolor{red}{Integrante 1}}
\addauthor{\textcolor{red}{Integrante 2}}
\addauthor{\textcolor{red}{Integrante 3}}

%% Forma de citar este trabalho: SOBRENOME, Nome (ou SOBRENOME et al., se mais de 3 autores)
\authorbib{\textcolor{red}{SOBRENOME, Nome et al.}}

%% Título (capa e folha de rosto) e título em inglês
\title{Propagação de fake news em sistemas paralelos e distribuídos: simulação por autômato celular com threads e sockets}
\titleen{Fake news propagation in parallel and distributed systems: cellular automaton simulation with threads and sockets}

%% Ano (capa e folha de rosto)
\depositshortdate{2026}

%% Instituição
\university{Universidade Tecnológica Federal do Paraná}
\universityen{Federal University of Technology -- Paraná}
\universitycampus{Campus Cornélio Procópio}
\universityunit{Departamento Acadêmico de Computação}
\address{Cornélio Procópio}
\addressen{Campus Cornélio Procópio, PR, Brazil}
\documenttype{Relatório Técnico}
\documenttypeen{Technical Report}
\degreetype{Graduação}

%% Palavras-chave (resumo) e keywords (abstract): informe a quantidade (máximo 5);
%% a classe cuida da pontuação. Use os comandos de acento do LaTeX (\'a, \^o, \~a, \c{c}).
\NumeroDePalavrasChave{4}
\PalavraChaveA{aut\^omatos celulares}
\PalavraChaveB{computa\c{c}\~ao paralela}
\PalavraChaveC{sistemas distribu\'idos}
\PalavraChaveD{propaga\c{c}\~ao de desinforma\c{c}\~ao}
%\PalavraChaveE{...}

\NumeroDeKeywords{4}
\KeywordA{cellular automata}
\KeywordB{parallel computing}
\KeywordC{distributed systems}
\KeywordD{misinformation spread}
%\KeywordE{...}

%% Licença Creative Commons: exigida só em TCC entregue à biblioteca. Vazio = sem licença.
%% Opções: ccby, ccbysa, ccbynd, ccbync, ccbyncsa, ccbyncnd (imagens em PreTexto/arquivos/licencasCC).
\licenca{}

